%=====================================================================
\chapter{The Semester Project}\label{ch:project}
%=====================================================================
\locator{6}{Part VI --- Emerging Hardware and Practice}

\begin{outcomes}
By the end of this chapter you will be able to:
\begin{tight}
  \item \textbf{Choose} a project that can be built, measured and defended in
        the time available.
  \item \textbf{State} a hypothesis your deployment will test, with the
        measurement that settles it.
  \item \textbf{Build} a virtualized service and \textbf{measure} it honestly,
        controlling for the five ways a virtualized benchmark lies.
  \item \textbf{Write} a report whose conclusions are supported by your own
        figures.
  \item \textbf{Present} and \textbf{defend} the work, including what did not
        work.
\end{tight}
\end{outcomes}

\begin{prereq}
The whole book. The project is assessed on whether you can apply it, and the
strongest projects are the ones that use the arithmetic of \chref{cpumem},
\chref{storage} and \chref{performance} on something the student built
themselves.
\end{prereq}

\begin{keyterms}
hypothesis $\bullet$ baseline $\bullet$ control $\bullet$ confounder $\bullet$
repeatability $\bullet$ percentile $\bullet$ teardown $\bullet$ reproducibility
$\bullet$ scope $\bullet$ deliverable $\bullet$ rubric $\bullet$ peer review
\end{keyterms}

\section{What the Project Is}

Build a virtualized or containerized service, measure something about it that
you did not know before you started, and report what you found.

Three deliverables, listed in the order they are worth marks:

\begin{tight}
  \item a \textbf{working deployment}, demonstrated live in Week 16;
  \item a \textbf{report} of 2{,}500--3{,}500 words containing your own
        measurements;
  \item a \textbf{ten-minute presentation} with five minutes of questions.
\end{tight}

The project is set in Week 9 and due in Week 16 (Appendix~G). Teams of two or
three; individual work is permitted but not encouraged, because the defence is
easier when somebody has already argued with you about it.

\begin{alertbox}[The commonest way to lose marks is scope]
The strongest projects in this course are narrow. ``Measure the throughput cost
of three virtio queue configurations'' earns more than ``build a private
cloud'', because the first can be finished, measured and defended in seven
weeks and the second cannot.

\smallskip
A project that builds a great deal and measures nothing is the single most
common failure mode. The deployment is the apparatus; the measurement is the
work.
\end{alertbox}

\section{Choosing One}

A good project has a \term{hypothesis} --- a statement that could turn out to be
false --- and a measurement that settles it. Here are examples at the right
scale; your own, if it has the same shape, is better.

\rowcolors{2}{white}{lightbg}
\begin{longtable}{@{}L{4.4cm}L{5.4cm}L{4.5cm}@{}}
\toprule
\rowcolor{primary!15}
\textbf{Project} & \textbf{The hypothesis it tests} & \textbf{Chapters} \\
\midrule
\endhead
Device model throughput & Paravirtualized devices cost less host CPU per unit
of work than emulated ones, and the ratio depends on request size &
\chref{paravirt}, \chref{performance} \\
Container density ceiling & Memory, not CPU, limits container density, and the
ratio to VMs depends on the workload's own footprint & \chref{containers} \\
Live migration convergence & Migration succeeds iff the dirty rate is below the
transfer rate, and auto-convergence buys success with guest performance &
\chref{migration} \\
Thin pool behaviour at exhaustion & Guest filesystems fail together, not
individually, and \texttt{discard} measurably changes the time to exhaustion &
\chref{storage} \\
Overlay MTU and throughput & The cost of each of the three MTU fixes, measured,
including the one that only helps TCP & \chref{network} \\
Isolation overhead & gVisor's cost is proportional to system-call rate, not to
CPU work & \chref{lightweight} \\
Right-sizing an estate & A measured estate needs materially less than it was
allocated, and \%RDY falls when vCPUs are removed & \chref{cpumem},
\chref{performance} \\
Cloud cost model & The three-year TCO of a real workload, on premises and in
two clouds, with the assumption that dominates identified & \chref{cloud} \\
\bottomrule
\end{longtable}

\begin{conceptbox}[The test for a hypothesis]
Write your hypothesis down and then ask: \emph{what result would prove this
wrong?}

\smallskip
If you cannot answer, it is not a hypothesis --- it is a plan to build
something. ``Containers are lighter than VMs'' fails the test, because no
measurement could contradict it. ``On this workload, the container-to-VM memory
density ratio is under 4:1 because the application's own footprint dominates''
passes: it names a number, and the number could come out at 20:1 and refute
you.

\smallskip
A refuted hypothesis, honestly reported, earns full marks. A confirmed one you
could not have refuted earns few.
\end{conceptbox}

\section{Measuring Honestly}

This is where most of the marks are, and \chref{performance} is the chapter to
re-read. Five controls, and your report must state how you handled each.

\rowcolors{2}{white}{lightbg}
\begin{longtable}{@{}L{3.2cm}L{5.1cm}L{6.0cm}@{}}
\toprule
\rowcolor{primary!15}
\textbf{Threat} & \textbf{What it does to your result} &
\textbf{What to do about it} \\
\midrule
\endhead
Neighbours & The host is shared, so your figure includes everyone else's work &
Run at several times of day; report the distribution; state what else was
running \\
Timekeeping & A descheduled guest's own clock drifts & Time from the host where
possible; use a paravirtual clock source \\
Caching & The host page cache serves your ``disk'' reads & Use files larger than
host RAM, or \texttt{direct=1}; state which \\
Burst credits & A short cloud benchmark runs on credits & Never use burstable
instances; say which family you used \\
Nesting & The outer hypervisor inflates everything, variably & Measure on
bare metal, or state clearly that figures are nested and comparative only \\
\bottomrule
\end{longtable}

\begin{conceptbox}[Three rules that separate a good report from a weak one]
\textbf{Report percentiles, not means.} \chref{performance} showed why: the
mean sits below almost all the mass. Give p50, p95 and p99, and say how many
samples.

\smallskip
\textbf{Report the baseline.} A number with nothing to compare it to is not a
measurement. Every result should be \emph{against native}, \emph{against the
default configuration}, or \emph{against the other arm}.

\smallskip
\textbf{Repeat, and show the spread.} Three runs minimum. If your three runs
differ by more than the effect you are claiming, you have not measured the
effect --- you have measured the noise, and saying so is a legitimate finding.
\end{conceptbox}

\section{The Report}

Roughly this shape. Marks follow the rubric below, not the word count.

\begin{enumerate}[leftmargin=2.2em, itemsep=3pt, topsep=4pt]
  \item \textbf{The question} (200 words). The hypothesis, and why it is worth
        testing. Name the chapter it comes from.
  \item \textbf{Background} (400). What the reader needs from this book to
        follow you. Cite chapters rather than re-explaining them.
  \item \textbf{Method} (600). The apparatus, in enough detail that somebody
        else could rebuild it. Hardware, versions, configuration. The five
        controls above, each addressed explicitly.
  \item \textbf{Results} (700). Your figures. Tables and plots with units and
        sample counts. No interpretation in this section.
  \item \textbf{Discussion} (800). What the figures mean. Whether the
        hypothesis survived. What surprised you. What you would do differently.
  \item \textbf{Threats to validity} (300). What could make your conclusion
        wrong. Every honest report has this section and it costs no marks to
        write --- omitting it costs several.
  \item \textbf{Conclusion and further work} (200).
  \item \textbf{Appendix}: your scripts, your raw data, your teardown
        procedure.
\end{enumerate}

\begin{alertbox}[Negative and inconclusive results are full-mark results]
If your measurement shows no difference where you expected one, report that. If
your three runs disagreed and you ran out of time to find out why, report that
too, with what you would do next.

\smallskip
What loses marks is a conclusion the figures do not support --- claiming a 12\%
improvement from runs that varied by 30\%, or reporting a mean that hides the
tail. The examiners can read your appendix.
\end{alertbox}

\section{The Presentation}

Ten minutes, five for questions. Expect to be asked: \emph{how do you know?},
\emph{what else could explain that?}, and \emph{what would change your
conclusion?}

Peer review is part of Week 16: each team reviews another's deployment against
the rubric and submits half a page. It is marked on the quality of the review,
not on how generous it is.

\begin{conceptbox}[Demonstrate live, and have a recording]
The deployment is demonstrated running. Things fail in front of audiences, so
record a working run beforehand and have it ready --- using it costs nothing,
and having nothing to show costs a great deal.

\smallskip
If something breaks live, say what you think happened. Diagnosing a failure in
front of the room, out loud, using the counters from \chref{performance}, is
the single most convincing thing that can happen in a project presentation.
\end{conceptbox}

\section{Marking}

\rowcolors{2}{white}{lightbg}
\begin{longtable}{@{}L{3.6cm}L{1.6cm}L{9.5cm}@{}}
\toprule
\rowcolor{primary!15}
\textbf{Component} & \textbf{Weight} & \textbf{What earns the marks} \\
\midrule
\endhead
Deployment & 25\% & It works, it is reproducible from your own instructions,
and it is torn down cleanly. Complexity earns nothing by itself \\
Measurement & 30\% & A real hypothesis; controls addressed; percentiles;
baseline; repetition. \textbf{The largest single component} \\
Report & 25\% & Clear, structured, conclusions supported by the figures, honest
threats-to-validity section \\
Presentation and defence & 15\% & Explained clearly; questions answered;
failures discussed rather than hidden \\
Peer review & 5\% & A specific, useful, fair review of another team \\
\bottomrule
\end{longtable}

\begin{conceptbox}[What distinguishes a first from a pass]
A passing project builds something that works and reports what it did.

\smallskip
A first-class project does three further things. It states a hypothesis that
could have been refuted. It controls for at least three of the five threats
and says so. And it reaches a conclusion the author did not expect at the
start --- which is only possible if the measurement was capable of surprising
them.

\smallskip
Notice that none of the three requires more building. They require deciding, at
the outset, what you are trying to find out.
\end{conceptbox}

\begin{pitfall}
\begin{tight}
  \item \textbf{Building for six weeks and measuring in the last one.} Build
        for two, measure for four. The measurement is where the marks are.
  \item \textbf{A scope that cannot be finished.} ``Build a private cloud'' is
        not a project; one measured property of one is.
  \item \textbf{Measuring once.} One run is an anecdote.
  \item \textbf{Reporting means.} Percentiles, always.
  \item \textbf{No baseline.} A figure with nothing to compare it to says
        nothing.
  \item \textbf{Benchmarking in a nested guest without saying so.} State it,
        and treat the figures as comparative only.
  \item \textbf{Hiding what failed.} The failure is usually the most
        interesting part, and the examiners will ask.
  \item \textbf{Leaving cloud resources running.} Every lab in this book ends
        with a teardown step for a reason. A forgotten instance bills until
        somebody notices.
\end{tight}
\end{pitfall}

\begin{lab}[Week 9: Your Proposal]
One page, submitted in Week 9, and the thing most likely to determine your
final mark. Expect ninety minutes, and do it with your team.

\begin{lstlisting}[language=bash]
# ch23_proposal.md -- one page, submitted in Week 9.
# Fill in every section. If you cannot fill one in, that is the
# section to work on, not the one to delete.

# --- 1. The hypothesis --------------------------------------------
# One sentence, containing a number or a comparison, that could turn
# out to be FALSE.
#   Bad : "Containers are more efficient than virtual machines."
#   Good: "On a 256 MB Flask service, container-to-VM memory density
#          is below 4:1, because the application footprint dominates
#          the per-instance overhead."

# --- 2. What would refute it --------------------------------------
# State the result that would prove you wrong. If you cannot, go
# back to 1.

# --- 3. The apparatus ---------------------------------------------
# Hardware you will use. Hypervisor and versions. How many guests.
# What you will install. Say explicitly whether it is nested.

# --- 4. The measurement -------------------------------------------
# What exactly you will measure, with which tool, how many times,
# and which percentiles you will report.
#   metric:      ...
#   tool:        ...
#   repetitions: at least 3
#   report:      p50, p95, p99 and the spread

# --- 5. The controls ----------------------------------------------
# One line each on the five threats of section 23.3:
#   neighbours:   ...
#   timekeeping:  ...
#   caching:      ...
#   burst credits:...
#   nesting:      ...

# --- 6. The baseline ----------------------------------------------
# What are you comparing against? Native? Default configuration?
# The other arm?

# --- 7. Weekly plan -----------------------------------------------
#   wk 10  proposal agreed; apparatus specified
#   wk 11  deployment built; first measurement attempted
#   wk 12  measurement working; baseline captured   <-- gate
#   wk 13  full runs; experiments FROZEN at the end of this week
#   wk 14  analysis; report drafted
#   wk 15  report finished; presentation rehearsed
#   wk 16  demonstrate, present, peer review
# The week 13 freeze is not a suggestion. Teams that are still
# changing the apparatus in week 15 have no time to analyse it.

# --- 8. What could go wrong ---------------------------------------
# Name the two most likely failures and what you will do instead.
# "The GPU is unavailable" -> "use the captured output in code/".

# --- 9. Teardown --------------------------------------------------
# How every resource you create will be destroyed, and when.
\end{lstlisting}

\textbf{What to notice.} Sections 1, 2 and 5 are the ones that decide the
mark, and they are the ones with no building in them. A proposal that is strong
on 3 and weak on 1 describes a deployment, not a project.
\end{lab}

\begin{checkpoint}
\begin{tightnum}
  \item Rewrite this as a testable hypothesis: ``We will investigate whether
        Kubernetes is good for our department.''
  \item A team reports a 9\% throughput improvement from a configuration
        change. Their three runs gave 7\%, 24\% and $-2$\%. What should the
        report conclude, and what should they do next?
  \item Your project needs a GPU and none is available in Week 12. Using the
        proposal's section 8, write what you would do instead, without changing
        the hypothesis.
\end{tightnum}
\end{checkpoint}

\begin{chaptersummary}
\begin{tight}
  \item Build a service, measure something you did not know, report it. The
        measurement is 30\% of the mark and the deployment 25\%.
  \item A hypothesis is a statement that could be refuted. If no result would
        prove you wrong, you have a build plan, not a project.
  \item Narrow scope beats ambition. Build for two weeks, measure for four.
  \item Control for neighbours, timekeeping, caching, burst credits and
        nesting, and say in the report how you did.
  \item Percentiles, a baseline, and at least three repetitions. If the runs
        vary more than the effect, you measured noise --- and saying so is a
        result.
  \item Negative and inconclusive results earn full marks. Conclusions the
        figures do not support do not.
  \item Freeze the apparatus in Week 13. Teams still changing it in Week 15
        have no time to analyse anything.
\end{tight}
\end{chaptersummary}

\begin{reviewq}
  \item \textbf{(MCQ)} The largest single component of the project mark is:
        (a)~the deployment; (b)~the measurement; (c)~the report;
        (d)~the presentation.
  \item \textbf{(MCQ)} A hypothesis is testable if: (a)~it is about a popular
        technology; (b)~some possible result would prove it false; (c)~it can
        be built in seven weeks; (d)~it cites three chapters.
  \item \textbf{(Short)} State the five threats to a virtualized measurement
        and one control for each.
  \item \textbf{(Short)} Explain why a refuted hypothesis can earn full marks
        and a confirmed one may not.
  \item \textbf{(Short)} Why is the Week 13 experiment freeze in the plan?
  \item \textbf{(Applied)} Take any one project from the table in Section 23.2
        and write its proposal sections 1, 2, 4 and 6 in full.
  \item \textbf{(Applied)} You are peer-reviewing a team whose report claims
        containers start 300 times faster than virtual machines, from one
        measurement of each, taken on a shared laboratory machine during a
        lecture. Write the half-page review: what is sound, what is not, and
        the two changes that would most improve it.
\end{reviewq}
